\subsection{Simple Regression Using Displacement}

Simple linear regression was performed in both Python using
\texttt{statsmodels} and Scala using ScalaTion. The two implementations
produced essentially identical results. The fitted equation was

\[
\widehat{MPG}
=
35.1206 - 0.0601(\mathrm{Displacement}).
\]

The negative slope indicates that greater engine displacement is associated
with lower fuel economy. A one-unit increase in displacement is associated
with an estimated decrease of approximately 0.0601 MPG. The slope was
statistically significant ($p<0.001$).

ScalaTion reported $R^2=0.648229$ and adjusted $R^2=0.647327$, consistent
with the statsmodels results. Thus, displacement alone explained approximately
64.8\% of the observed variation in MPG. ScalaTion also reported an RMSE of
approximately 4.62 MPG.

\subsection{Simple Regression Using Weight}

The simple regression using vehicle weight also produced nearly identical
results in ScalaTion and statsmodels. The fitted equation was

\[
\widehat{MPG}
=
46.2165 - 0.007647(\mathrm{Weight}).
\]

The negative slope indicates that heavier vehicles tend to have lower fuel
economy. A one-unit increase in vehicle weight is associated with an estimated
decrease of approximately 0.00765 MPG. The slope was statistically significant
($p<0.001$).

ScalaTion reported $R^2=0.692630$ and adjusted $R^2=0.691842$, again matching
the statsmodels results. Therefore, vehicle weight alone explained
approximately 69.3\% of the variation in MPG. ScalaTion reported an RMSE of
approximately 4.32 MPG.

\subsection{Regression Discussion}

Both selected predictors exhibited strong and statistically significant
negative relationships with MPG, and the ScalaTion and statsmodels
implementations produced essentially identical parameter estimates and
goodness-of-fit measures. Vehicle weight provided the stronger of the two
simple regression models, with $R^2 \approx 0.693$, compared with
$R^2 \approx 0.648$ for displacement. This result is consistent with the
correlation analysis, in which weight also had the stronger absolute
correlation with MPG. These results suggest that vehicle weight is the
stronger individual linear predictor of fuel economy among the two selected
variables.